\usepackage{booktabs}
\usepackage{array}
\usepackage{longtable}
\usepackage{tabularx}
\usepackage{amsmath}
\usepackage{amssymb}
\usepackage{pgfplots}
\pgfplotsset{compat=1.18}
\usepackage{tcolorbox}
\tcbuselibrary{skins,breakable}
\usepackage{hyperref}
\hypersetup{
    colorlinks=true,
    linkcolor=blueColor,
    urlcolor=blueColor,
    citecolor=purpleColor
}


\documentclass[12pt]{article}
\usepackage{amssymb}
\usepackage{geometry}
\usepackage{graphicx}
\usepackage[colorlinks=true, allcolors=blue]{hyperref}
\usepackage{cite}
\usepackage{wrapfig}
\usepackage{booktabs}
\usepackage{soul}
\usepackage[most]{tcolorbox}
\usepackage{amsmath}
\usepackage{lipsum}
\usepackage{caption}
\usepackage{ulem}
\usepackage{xcolor}
\usepackage{tikz}
\usetikzlibrary{arrows.meta, positioning, shapes}
\usepackage{sectsty}
\usepackage{booktabs}
\usepackage{array}
\usepackage{longtable}
\usepackage{tabularx}
\usepackage{amsmath}
\usepackage{amssymb}
\usepackage{pgfplots}
\pgfplotsset{compat=1.18}
\usepackage{tcolorbox}
\tcbuselibrary{skins,breakable}
\usepackage{hyperref}
\hypersetup{
    colorlinks=true,
    linkcolor=blueColor,
    urlcolor=blueColor,
    citecolor=purpleColor
}

%--------- Define Colors ---------
\definecolor{titleColor}{HTML}{800000}        % Maroon
\definecolor{sectionColor}{HTML}{4682B4}      % SteelBlue
\definecolor{subsectionColor}{HTML}{6A5ACD}   % SlateBlue
\definecolor{textHighlight}{HTML}{DAA520}     % Goldenrod
\definecolor{backgroundColor}{HTML}{F0F8FF}   % AliceBlue
\definecolor{emphasisColor}{RGB}{200, 50, 50} % ForestGreen

\sectionfont{\color{sectionColor}}       % Color for \section
\subsectionfont{\color{subsectionColor}} % Color for \subsection

%--------- Custom Highlight Commands ---------
\definecolor{lightblue}{RGB}{230, 240, 255}
\definecolor{lightgreen}{RGB}{220, 255, 220}
\definecolor{lightred}{RGB}{255, 230, 230}
\definecolor{lightyellow}{RGB}{255, 255, 200}
\definecolor{darkblue}{RGB}{0, 51, 102}

\newcommand{\bluehighlight}[1]{\colorbox{lightblue}{#1}}
\newcommand{\greenhighlight}[1]{\colorbox{lightgreen}{#1}}
\newcommand{\redhighlight}[1]{\colorbox{lightred}{#1}}
\newcommand{\yellowhighlight}[1]{\colorbox{lightyellow}{#1}}
\newcommand{\blue}[1]{\textcolor{darkblue}{#1}}
\newcommand{\myemphasis}[1]{\textcolor{emphasisColor}{\textbf{#1}}}

%--------- Color Blocks ---------
\newtcolorbox{blockA}{
  breakable,
  colback=blue!5,
  colframe=blue!40!black,
  enhanced,
  boxrule=0.7pt,
  arc=2mm,
  left=4mm,right=4mm,top=2mm,bottom=2mm
}

\newtcolorbox{blockB}{
  breakable,
  colback=red!5,
  colframe=red!40!black,
  enhanced,
  boxrule=0.7pt,
  arc=2mm,
  left=4mm,right=4mm,top=2mm,bottom=2mm
}

\newtcolorbox{blockC}{
  breakable,
  colback=green!5,
  colframe=green!40!black,
  enhanced,
  boxrule=0.7pt,
  arc=2mm,
  left=4mm,right=4mm,top=2mm,bottom=2mm
}

\newtcolorbox{blockD}{
  breakable,
  colback=purple!5,
  colframe=purple!40!black,
  enhanced,
  boxrule=0.7pt,
  arc=2mm,
  left=4mm,right=4mm,top=2mm,bottom=2mm
}

\newcommand{\ABCTotalMedicines}{50}
\newcommand{\ABCTotalValue}{1,474,610.34}
\newcommand{\ABCObservationDays}{1348}
\newcommand{\ABCACount}{16}
\newcommand{\ABCAValue}{1,157,967.33}
\newcommand{\ABCAShare}{78.53\%}
\newcommand{\ABCBCount}{16}
\newcommand{\ABCBValue}{241,111.57}
\newcommand{\ABCBShare}{16.35\%}
\newcommand{\ABCCCount}{18}
\newcommand{\ABCCValue}{75,531.44}
\newcommand{\ABCCShare}{5.12\%}


\sectionfont{\color{titleColor}}
\subsectionfont{\color{blueColor}}

\title{ \textcolor{titleColor}{ \Huge\bfseries REEM V2\\[0.3cm] ABC Inventory Intelligence }}

\author{ Research and Development Report\\ REEM Intelligence Engine }
\date{\today}

\begin{document}

\maketitle

\begin{tcolorbox}[
    colback=titleColor!5,
    colframe=titleColor,
    title=\textbf{Development Status},
    breakable
]
This report documents the ABC Inventory Intelligence component of REEM V2.
The underlying dataset is synthetic development data generated for model
development and validation. It is not operational pharmacy evidence and must
not be used for real procurement or inventory decisions.
\end{tcolorbox}

\newpage

\section{Executive Summary}

ABC analysis provides a value-based prioritisation of medicines according to
their observed consumption value. It complements the REEM Inventory
Intelligence model, which focuses on stock sufficiency and reorder decisions.

The current development dataset contains
\textbf{\ABCTotalMedicines\ medicines} observed over
\textbf{\ABCObservationDays\ days}. Annualised consumption value is estimated
from observed sales quantity and medicine unit cost.

The total annualised consumption value is
\textbf{\ABCTotalValue}.

The resulting classification is:

\begin{center}
\begin{tabular}{lrrr}
\toprule
Class & Medicines & Value & Share\\
\midrule
A & 16 & 1,157,967.33 & 78.52\%\\
B & 16 & 241,111.57 & 16.35\%\\
C & 18 & 75,531.44 & 5.12\%\\
\bottomrule
\end{tabular}
\end{center}

The A class contains only 16 medicines but represents 78.52\% of annualised consumption value.  This makes A-class medicines the primary management priority in a value-based inventory control framework.

\section{Purpose}

The purpose of ABC Inventory Intelligence is to transform pharmacy
consumption evidence into a management-priority structure.

The analytical question is:

\begin{quote}
\textit{Which medicines account for the largest share of consumption value
and therefore deserve the greatest management attention?}
\end{quote}

ABC analysis is not intended to replace demand forecasting, reorder-point
analysis, safety-stock modelling, expiry analysis, or supplier intelligence.
It is a complementary management layer.

\section{Position in the REEM Architecture}

\begin{center}
\begin{tikzpicture}[
    node distance=1.2cm,
    box/.style={
        rectangle,
        rounded corners,
        draw=blueColor,
        thick,
        minimum width=5.2cm,
        minimum height=0.9cm,
        align=center
    },
    arrow/.style={
        -{Latex[length=3mm]},
        thick,
        draw=purpleColor
    }
]

\node[box] (data) {Operational Pharmacy Data\\Sales and Consumption};
\node[box, below=of data] (quality) {Data Quality\\Observation Period};
\node[box, below=of quality] (value) {Consumption Value\\Quantity $\times$ Unit Cost};
\node[box, below=of value] (annual) {Annualisation\\365-Day Equivalent};
\node[box, below=of annual] (abc) {ABC Classification\\A / B / C};
\node[box, below=of abc] (decision) {Management Priority\\HIGH / MEDIUM / LOW};

\draw[arrow] (data) -- (quality);
\draw[arrow] (quality) -- (value);
\draw[arrow] (value) -- (annual);
\draw[arrow] (annual) -- (abc);
\draw[arrow] (abc) -- (decision);

\end{tikzpicture}
\end{center}

\section{Evidence Architecture}

ABC V1 uses observed sales as the consumption evidence.

For each medicine:

\begin{equation}
V_i = Q_i \times C_i
\end{equation}

where:

\begin{itemize}
    \item $Q_i$ is observed quantity consumed,
    \item $C_i$ is unit acquisition cost,
    \item $V_i$ is observed consumption value.
\end{itemize}

Because the observation period is not exactly one year, consumption is
annualised:

\begin{equation}
Q_{i,annual}
=
Q_i
\left(
\frac{365}{D}
\right)
\end{equation}

where $D$ is the observed number of calendar days.

Annualised consumption value is therefore:

\begin{equation}
V_{i,annual}
=
Q_{i,annual} \times C_i
\end{equation}

This produces a common annual-equivalent scale for ranking medicines.

\section{Why Consumption Value Is Used}

ABC analysis is fundamentally a value-prioritisation method.

The model deliberately does not classify medicines according to:

\begin{itemize}
    \item current stock quantity,
    \item current stock value,
    \item forecast demand alone,
    \item medicine criticality alone,
    \item selling price alone.
\end{itemize}

Instead, it combines observed consumption quantity with unit cost.

This distinction is important because a medicine may have high physical
volume but relatively low financial consumption, while another medicine may
have lower unit movement but substantially higher consumption value.

\section{Classification Method}

Medicines are sorted in descending order of annualised consumption value.

For each medicine, its percentage contribution is calculated as:

\begin{equation}
P_i =
\frac{V_{i,annual}}
{\sum_j V_{j,annual}}
\times 100
\end{equation}

The cumulative percentage is then calculated in descending value order.

The development thresholds are:

\begin{center}
\begin{tabular}{lll}
\toprule
Class & Cumulative value boundary & Priority\\
\midrule
A & up to approximately 80\% & HIGH\\
B & above A to approximately 95\% & MEDIUM\\
C & remaining value & LOW\\
\bottomrule
\end{tabular}
\end{center}

The implementation uses the cumulative threshold at each ranked medicine,
so the resulting class boundary is data-driven rather than forcing an exact
number of medicines into each class.

\section{Current Development Results}

The current REEM V2 synthetic world contains:

\begin{itemize}
    \item \textbf{50 medicines}
    \item \textbf{1348 days} of observation
    \item \textbf{1,474,610.34} annualised consumption value
    \item \textbf{16 A-class medicines}
    \item \textbf{16 B-class medicines}
    \item \textbf{18 C-class medicines}
\end{itemize}

The A class represents \textbf{78.52\%} of annualised consumption value.
The B class represents \textbf{16.35\%} of annualised consumption value.

The C class represents \textbf{5.12\%} of annualised consumption value.

\section{ABC Value Distribution}

\begin{center}
\begin{tikzpicture}
\begin{axis}[
    width=0.92\textwidth,
    height=8cm,
    xlabel={Medicine rank by annualised consumption value},
    ylabel={Cumulative consumption value (\%)},
    xmin=1,
    xmax=50,
    ymin=0,
    ymax=105,
    grid=major,
    thick
]
\addplot[
    mark=*,
    mark size=1.2pt
] coordinates {
(1,11.2500) (2,19.8700) (3,26.3100) (4,32.5200) (5,38.4400) (6,43.6900) (7,48.3800) (8,52.8000) (9,57.0500) (10,60.6600) (11,64.2500) (12,67.7800) (13,71.0100) (14,73.6300) (15,76.2400) (16,78.5300) (17,80.4100) (18,82.0900) (19,83.5200) (20,84.7600) (21,85.9300) (22,87.0400) (23,88.0600) (24,89.0700) (25,89.9100) (26,90.7200) (27,91.5300) (28,92.3000) (29,93.0500) (30,93.7300) (31,94.3100) (32,94.8800) (33,95.4300) (34,95.9800) (35,96.4100) (36,96.7900) (37,97.1700) (38,97.5200) (39,97.8500) (40,98.1700) (41,98.4500) (42,98.7200) (43,98.9600) (44,99.1700) (45,99.3600) (46,99.5300) (47,99.6900) (48,99.8500) (49,99.9900) (50,100.0000) 
};
\addplot[
    dashed,
    domain=1:50
] coordinates {(1,80) (50,80)};
\addplot[
    dashed,
    domain=1:50
] coordinates {(1,95) (50,95)};
\end{axis}
\end{tikzpicture}
\end{center}

The curve demonstrates the concentration of consumption value among the
highest-ranked medicines. The A boundary captures the dominant portion of
value, while the remaining medicines account for progressively smaller
individual contributions.

\section{Top Consumption-Value Medicines}

\begin{longtable}{rllrrrr}
\toprule
Rank & ID & Medicine & Annual Value & Share & Cumulative & Class\\
\midrule
1 & D0018 & Carbamazepine & 165,820.62 & 11.25\% & 11.25\% & A \\
2 & D0021 & Cetirizine & 127,136.19 & 8.62\% & 19.87\% & A \\
3 & D0016 & Budesonide & 95,021.44 & 6.44\% & 26.31\% & A \\
4 & D0048 & VitaminD & 91,523.33 & 6.21\% & 32.52\% & A \\
5 & D0043 & Salbutamol & 87,318.94 & 5.92\% & 38.44\% & A \\
6 & D0005 & Amikacin & 77,489.44 & 5.25\% & 43.69\% & A \\
7 & D0011 & Aspirin & 69,086.29 & 4.69\% & 48.38\% & A \\
8 & D0003 & Allopurinol & 65,152.50 & 4.42\% & 52.80\% & A \\
9 & D0007 & Amlodipine & 62,663.55 & 4.25\% & 57.05\% & A \\
10 & D0046 & Tramadol & 53,293.02 & 3.61\% & 60.66\% & A \\

\bottomrule
\end{longtable}

The ranking identifies the medicines that should receive the strongest
value-based management attention.

\section{Management Interpretation}

The current development results support a simple management principle:

\begin{tcolorbox}[
    colback=goldColor!8,
    colframe=goldColor,
    title=\textbf{Management Principle},
    breakable
]
Not every medicine requires the same intensity of inventory control.
Management effort should be concentrated according to the economic importance
of the medicine within the observed consumption portfolio.
\end{tcolorbox}

\subsection{A-Class Medicines}

A-class medicines represent the largest share of consumption value.

Recommended management attention includes:

\begin{itemize}
    \item closer monitoring of demand,
    \item stronger procurement planning,
    \item tighter stock-control thresholds,
    \item regular review of supplier performance,
    \item stronger expiry surveillance,
    \item more frequent management review.
\end{itemize}

\subsection{B-Class Medicines}

B-class medicines have intermediate economic importance.

They can normally receive standard inventory-control procedures while remaining
visible to management through periodic review.

\subsection{C-Class Medicines}

C-class medicines represent the smaller share of consumption value.

They may be managed using simpler control procedures, provided that clinical
criticality and other risk factors are considered separately.

\section{ABC Is Not Clinical Criticality}

ABC classification must not be interpreted as a measure of clinical
importance.

A low-value medicine can still be clinically essential.

Therefore:

\begin{equation}
\text{Management Priority}
\neq
\text{Clinical Criticality}
\end{equation}

REEM should ultimately combine ABC value priority with medicine criticality,
expiry risk, stockout risk, and demand uncertainty.

\section{Relationship With Inventory Intelligence}

ABC and Inventory Intelligence answer different questions.

\begin{center}
\begin{tikzpicture}[
    node distance=1.5cm,
    box/.style={
        rectangle,
        rounded corners,
        draw=blueColor,
        thick,
        minimum width=5.0cm,
        minimum height=1.0cm,
        align=center
    },
    arrow/.style={
        -{Latex[length=3mm]},
        thick,
        draw=purpleColor
    }
]

\node[box] (abc) {ABC Intelligence\\
Which medicines matter most by value?};

\node[box, right=2cm of abc] (inv) {Inventory Intelligence\\
Which medicines need stock action?};

\node[box, below=1.7cm of abc] (combined) {REEM Management Intelligence\\
What deserves attention and what action is appropriate?};

\draw[arrow] (abc) -- (combined);
\draw[arrow] (inv) -- (combined);

\end{tikzpicture}
\end{center}

This separation is central to the REEM architecture.

ABC provides prioritisation.

Inventory Intelligence provides stock-state interpretation.

Together they form a stronger management decision layer.

\section{Evidence Quality}

The current ABC model is analytically coherent, but its evidence status remains
developmental.

\begin{center}
\begin{tabular}{ll}
\toprule
Dimension & Status\\
\midrule
Data source & SIMULATED\_REEM\_V2\\
Data status & DEVELOPMENT\\
Observation period & 1348 days\\
Medicines analysed & 50\\
Classification completeness & 50/50\\
Operational use & PROHIBITED\\
\bottomrule
\end{tabular}
\end{center}

The classification is therefore a validated development result, not an
operational pharmacy finding.

\section{Model Integrity Validation}

The ABC validation confirmed:

\begin{itemize}
    \item 50 medicines were classified;
    \item A/B/C classes contain 16, 16, and 18 medicines respectively;
    \item total annualised value is 1,474,610.34;
    \item final cumulative value is 100\%;
    \item every medicine belongs to exactly one ABC class;
    \item the class shares sum to approximately 100\%.
\end{itemize}

These checks support the internal integrity of ABC V1.

\section{Limitations}

Several limitations remain.

\subsection{Synthetic Data}

The entire current dataset is synthetic and was generated for REEM
development. It does not represent real pharmacy operations.

\subsection{Observed Consumption}

The model assumes that observed sales provide a useful representation of
consumption.

Future operational versions should distinguish sales from true demand,
including:

\begin{itemize}
    \item unfilled demand,
    \item stockout periods,
    \item substitutions,
    \item returns,
    \item emergency requests,
    \item demand requests not converted into sales.
\end{itemize}

\subsection{Cost Assumption}

The current calculation uses medicine unit cost as the value basis.

Future versions may incorporate time-varying acquisition cost and procurement
conditions.

\subsection{Single-Dimension Classification}

ABC is intentionally a value-based classification. It should eventually be
combined with other dimensions rather than treated as a complete inventory
risk model.

\section{Next Development Stage}

The next REEM development stages can extend ABC toward a multidimensional
inventory intelligence framework:

\begin{enumerate}
    \item ABC value classification;
    \item criticality classification;
    \item demand variability;
    \item expiry risk;
    \item supplier performance;
    \item stockout risk;
    \item service-level requirements;
    \item integrated management priority.
\end{enumerate}

A future matrix could therefore distinguish:

\begin{center}
\begin{tabular}{lll}
\toprule
ABC Value & Clinical Criticality & Management Interpretation\\
\midrule
A & High & Highest attention\\
A & Low & High economic attention\\
C & High & High clinical attention\\
C & Low & Routine attention\\
\bottomrule
\end{tabular}
\end{center}

This would move REEM beyond conventional ABC analysis toward integrated
decision intelligence.

\section{Management Intelligence Perspective}

ABC analysis is valuable because it reduces a large inventory portfolio to a
manageable management structure.

Instead of treating 50 medicines equally, REEM identifies the smaller group
that accounts for most of the consumption value.

The analytical chain is:

\begin{center}
\begin{tikzpicture}[
    node distance=1.1cm,
    box/.style={
        rectangle,
        rounded corners,
        draw=purpleColor,
        thick,
        minimum width=4.5cm,
        minimum height=0.8cm,
        align=center
    },
    arrow/.style={
        -{Latex[length=3mm]},
        thick,
        draw=goldColor
    }
]

\node[box] (obs) {Observed Consumption};
\node[box, below=of obs] (val) {Consumption Value};
\node[box, below=of val] (rank) {Value Ranking};
\node[box, below=of rank] (class) {ABC Classification};
\node[box, below=of class] (priority) {Management Priority};

\draw[arrow] (obs) -- (val);
\draw[arrow] (val) -- (rank);
\draw[arrow] (rank) -- (class);
\draw[arrow] (class) -- (priority);

\end{tikzpicture}
\end{center}

This is the central contribution of ABC Intelligence within REEM V2.

\section{Conclusion}

ABC Inventory Intelligence V1 successfully transforms observed consumption
data into a value-based management classification.

In the current synthetic development world, 16 medicines form the A class and
represent approximately 78.5\% of annualised consumption value. Sixteen
medicines form the B class, while eighteen form the C class.

The model is internally validated and provides a clear management
prioritisation layer.

However, the evidence remains synthetic and developmental. The model must
therefore remain explicitly separated from operational pharmacy decision
making until it is tested against longitudinal real-world data.

The next stage is not simply to produce more ABC classifications. It is to
combine ABC value importance with demand, stock, expiry, supplier, and
clinical-risk evidence so that REEM can progress from classification toward
integrated management intelligence.

\vfill

\begin{tcolorbox}[
    colback=titleColor!5,
    colframe=titleColor,
    title=\textbf{REEM V2 Evidence Boundary},
    breakable
]
\centering
\textbf{SIMULATED\_REEM\_V2}\\
\textbf{DEVELOPMENT}\\
\textbf{OPERATIONAL USE PROHIBITED}
\end{tcolorbox}

\end{document}






\section{Current Development Results}

The current REEM V2 synthetic world contains:

\begin{itemize}
    \item \textbf{50 medicines}
    \item \textbf{1348 days} of observation
    \item \textbf{1,474,610.34} annualised consumption value
    \item \textbf{16 A-class medicines}
    \item \textbf{16 B-class medicines}
    \item \textbf{18 C-class medicines}
\end{itemize}

The A class contains only \textbf{16 medicines} but represents
\textbf{78.52\%} of annualised consumption value.

The B class represents \textbf{16.35\%} of annualised consumption value.

The C class represents \textbf{5.12\%} of annualised consumption value.

